% s03 The inputs
\section{The inputs}\label{sec:s03}

This section states, in the form in which the later sections consume them, the results on which the theorems of this paper rest, and says where each is proved. Three are theorems of this paper proved in later parts: the existence of King's cube limits, without and with scalar sources (\cref{hyp:s03:EC,hyp:s03:W1}; the existence part, \cref{sec:k01} and \moving{existence section with scalar sources}), and the rotation theorem (\cref{hyp:s03:RC}; the rotation part, \moving{rotation part}). Two are Ba{\l}aban's torus estimates in the form used by the Wick leg (\cref{hyp:s03:T1,hyp:s03:T2}). One is a condition on the couplings, the sign clause \Tthree\ (\cref{hyp:s03:T3}), proved on explicit coupling regions in \cref{prop:s14:T3}; it is the only condition on the couplings anywhere in the paper. The hypotheses of the optional extension of the Wick leg to arbitrary amplitudes (\Wtwo, \RW, \FAK) are stated at the end of the section; no clause of the main theorem uses them. From the literature, the paper uses the correlation inequalities of \cite{BFS1}, the chessboard estimate of \cite{FILS78}, the reconstruction theorems of \cite{OS73,OS75}, Ba{\l}aban's torus analysis \cite{Ba1,Ba2,Ba3,Ba4} and King's papers \cite{K1,K2}, each at the place where it is applied; the general lemmas on the measure, its ergodic decomposition and weak clustering are taken from the authors' earlier paper on the torus limits \cite{SymPhase} (\cref{rem:s03:companion-inputs}).

\subsection{Objects and conventions}

The two cube conventions of Definitions~\ref{def:s02:Kcube} and~\ref{def:s02:acube} are King's convention $\NDK$ (the scalar field killed on the three top face layers $\TJ$ of the cube $\CJ$, the convention of the state King prints) and the comparison convention $\NDa$ (the scalar alive on all of $\XJ$), in which the tube results (Theorem~\ref{thm:s11:W}, Lemma~\ref{lem:s11:TWtube}) and the lattice comparisons of the rotation part are written. The two are interlaced by Lemma~\ref{lem:s04:G}. Which convention a statement is posed on is part of it:
\begin{itemize}[leftmargin=2em]
\item The existence theorems \cref{hyp:s03:EC,hyp:s03:W1} are posed on King's convention $\NDK$, on which King's renormalisation-group flow runs (\cref{rem:k01:cube}). By Corollary~\ref{cor:s04:G2} the existence of the fixed-cube limits on one convention does not imply it on the other; the comparison family's fixed-cube limits are reached through that corollary, and the statements proved for the comparison family under the assumption that they exist (\cref{def:s03:conv}) serve as templates for King's convention.
\item The rotation theorem \cref{hyp:s03:RC} is proved for the comparison family $\NDa$ and reaches King's printed state $\ZKinf$ through Corollary~\ref{cor:s04:G2}.
\item \Cref{hyp:s03:T1,hyp:s03:T2} and \Tthree\ are statements about King's tori (Definition~\ref{def:s10:torus}) and about $c$-numbers.
\end{itemize}
Throughout, $Z^{a}_{k,J}(z,w)$ denotes the generating functional of Definition~\ref{def:s02:sources} with curvature coefficients $z$ and Wick coefficients $w$ in the scheme $a$, on the convention in question (written $Z^{a,K}_{k,J}$ when the convention $\NDK$ is to be stressed); $z=a^\sharp$ is the reference scheme of Remark~\ref{rem:s02:scheme}; $\avK{\,\cdot\,}$ is King's cube expectation at spacing $\eps_k$.

\begin{definition}[convergence of the comparison family at fixed cubes]\label{def:s03:conv}
The comparison family is \emph{convergent at fixed cubes} if for every finite family $f_1,\dots,f_n$ of real $2$-forms in $C^\infty_c(\R^3;\Lambda^2)$, every $J\ge J_0(f)$ and every $x\in\R^n$ the limit $\lim_k\ava{e^{\sum_jx_jF(f_j)}}$ exists and lies in $(0,\infty)$.
\end{definition}
This is the clause of \cref{hyp:s03:EC} read for the family $\NDa$. It is not a result of this paper and is not needed for the main theorem: the statements below that assume it are proved for the comparison family as templates, and their forms for King's printed state are proved under \cref{hyp:s03:EC}; the continuum-subsequence section (\cref{thm:s12:LIM}) removes the dependence on the fixed-cube limits altogether.

\subsection{Existence of King's cube limits}

\begin{theorem}[King's cube limits, curvature]\label{hyp:s03:EC}
For every finite family $f_1,\dots,f_n$ of real $2$-forms in $C^\infty_c(\R^3;\Lambda^2)$, every $J\ge J_0(f)$, where $J_0(f)$ is the least $J$ such that every $\supp f_j$ lies at $\ell^\infty$-distance greater than $1$ from $\R^3\setminus\CJ$, and every $z\in\R^n$, the limit
\[
\lim_{k\to\infty}\avK{e^{\sum_j z_j\Fcurv{f_j}}}
\]
exists and lies in $(0,\infty)$.
\end{theorem}
Proved in \cref{sec:k01} (\cref{thm:k01:EC}), from King's renormalisation-group analysis of the cube.

\begin{remark}[on \cref{hyp:s03:EC}]\label{rem:s03:EC}
(i) \emph{Existence clause only.} \Cref{hyp:s03:EC} is used only through the existence of the fixed-cube limit at real coefficients: no constants clause, no uniformity in $J$ and no locally uniform convergence in $z$ enter from it. Every bound used on the limit is proved here, with absolute constants, in Theorem~\ref{thm:s05:C}; the monotonicity in the cube, the $J$-limit, the entire extension and the translations are proved in Lemma~\ref{lem:s04:G}, Theorem~\ref{thm:s05:C}, Lemma~\ref{lem:s05:TWC} and Theorem~\ref{thm:s06:O3}.
(ii) \emph{Origin.} This is the existence clause of King's Theorem~2.1 \cite{K2} at zero scalar source, for the state he prints. The proof in \cref{sec:k01} follows King's route, the comparison of two spacings at a common block scale, in a ratio form in which the vacuum energies cancel; what it uses of Ba{\l}aban's block-factor-$5$ analysis \cite{Ba1,Ba2,Ba3,Ba4} is stated there. The mass clause of \cite[Prop.~3.1]{Ba2}, upstream of it, is part of that proposition's statement and is cited under it (Remark~\ref{rem:s14:BALm}); it is not a condition of this paper.
(iii) \emph{Consumers.} The whole curvature sector (Theorems~\ref{thm:s05:C}, \ref{thm:s06:O3}, \ref{thm:s08:OS2C}, \ref{thm:s09:S2K}), the identification of the LIM states with King's state (Corollary~\ref{cor:s12:K}) and the main theorem (Theorem~\ref{thm:s13:main}).
\end{remark}

\begin{remark}[the cube]\label{rem:s03:cube-identification}
The cube on which the existence theorem is proved is King's RG cube $\CK$ of \cref{def:s02:Kcube}, with King's killing of the scalar on the top layers $\TJ$; it is the only cube on which King's block-spin flow exists (\cref{rem:k01:cube}). Every placement of a comparison-cube result on King's cube goes through Lemma~\ref{lem:s04:G} and its corollaries.
\end{remark}

\begin{theorem}[King's cube limits with scalar sources]\label{hyp:s03:W1}
Let $z=a^\sharp$ be the reference scheme and let $s_W\defeq s_1$ be the amplitude of Proposition~\ref{prop:s10:U}. For every finite family of real $2$-forms $f_j\in C^\infty_c(\R^3;\Lambda^2)$ and real functions $h_i\in C^\infty_c(\R^3)$, every $J\ge J_0(f,h)$ and every real coefficient vector $(z,w)$ whose positive scalar part satisfies $\sup\big(\sum_iw_ih_i\big)_+\le s_W$ (the negative scalar part is unrestricted), the limit $\lim_{k\to\infty}Z^{z}_{k,J}(z,w)$ exists and lies in $(0,\infty)$.
\end{theorem}
Proved in \moving{existence section with scalar sources}, by the route of \cref{sec:k01} with the scheme $c$-number carried along.

\begin{definition}[admissible schemes]\label{def:s03:adm}
A scheme $a$ (Remark~\ref{rem:s02:scheme}) is \emph{admissible} if $a=z$ or $\lim_k(z_k-a_k)$ exists. For an admissible scheme \cref{hyp:s03:W1} holds for $Z^a_{k,J}$ in place of $Z^z_{k,J}$ (Lemma~\ref{lem:s02:alg}(a)), and $\sup_k|z_k-a_k|<\infty$.
\end{definition}

\begin{remark}[on \cref{hyp:s03:W1}]\label{rem:s03:W1}
(i) At zero Wick coefficient \cref{hyp:s03:W1} is \cref{hyp:s03:EC}.
(ii) \emph{Origin.} \Cref{hyp:s03:W1} is \emph{not} King's Theorem~2.1 with scalar sources: as printed, in King's Wick ordering, that statement fails, and the state with scalar sources exists only after a vacuum counterterm, which is what the scheme supplies (Remark~\ref{rem:s02:scheme}). The threshold on the positive scalar part is the one at which the positive-source bound of the Wick leg is available (Proposition~\ref{prop:s11:W1prime}); what happens above it is the subject of the optional extension (\cref{hyp:s03:W2}).
(iii) One instance of \cref{hyp:s03:W1} is named separately because the statements on continuum subsequences need it alone: \Bcell\ (Hypothesis~\ref{hyp:s12:Bcell}; one cube, one negative Wick source).
(iv) \emph{Consumers.} The Wick tube (Proposition~\ref{prop:s10:U}, Theorem~\ref{thm:s11:W}), its transfer to King's convention (Theorem~\ref{thm:s13:TWK}), the Wick sector of Theorem~\ref{thm:s12:LIM}, and Form~(B) of the main theorem (Theorem~\ref{thm:s13:main}, Proposition~\ref{prop:s13:WK}).
\end{remark}

\subsection{The torus inputs of the Wick leg}

\begin{proposition}[torus stability at two scalar masses]\label{hyp:s03:T1}
For the scalar masses $M=m$ and $M=M_1$ (the auxiliary mass fixed in Proposition~\ref{prop:s10:U}) there are constants $E_\pm(M)<\infty$, a tail $K_1$ and, for each $k\ge K_1$, a $J_*(k)$ such that the source-free partition functions of King's tori $T$ (Definition~\ref{def:s10:torus}) at spacing $\eps_k$ and at scalar mass $M$ satisfy, in King's torus convention,
\[
e^{-E_-(M)\abs{T}}\le Z^{(M)}_{k,T}\le e^{E_+(M)\abs{T}}\qquad(k\ge K_1,\ J\ge J_*(k)).
\]
\end{proposition}
These are Ba{\l}aban's lower and upper stability bounds, \cite[(1.14)]{Ba1} and \cite[(1.3)]{Ba2}, at a fixed scalar mass, read at the two masses $m$ and $M_1$ and in King's torus normalisation of the counterterm; the form used here is established in \moving{existence part: Ba{\l}aban's torus bounds at block factor $5$}. Proposition~3.1 of \cite{Ba2}, behind the upper bound, is stated under a cap on the dimensionless scalar mass at the scales at which it is applied; it is cited under that cap, one more one-sided cap on the stopping scale and no condition on the bare mass (Remark~\ref{rem:s14:BALm}). Only a tail in $k$ and, for each $k$, large $J$ are used (Proposition~\ref{prop:s10:U}, Theorem~\ref{thm:s11:W}).

\begin{lemma}[the $c$-number input]\label{hyp:s03:T2}
Let $\alpha^\infty_k(M)=M^2+\dm_k(M)$ be King's free-boundary bare coefficient at scalar mass $M$ and spacing $\eps_k$, in the reference scheme $z$. For the auxiliary mass $M_1$ of Proposition~\ref{prop:s10:U}, the quantities
\[
\sigma_1(M_1)\defeq\liminf_{k\to\infty}\tfrac12\bigl(\alpha^\infty_k(m)-\alpha^\infty_k(M_1)\bigr)
\qquad\text{and}\qquad \beta_1(M_1)
\]
are finite, where $\beta_1(M_1)$ is the second constant entering Proposition~\ref{prop:s10:U}.
\end{lemma}
Proved in \moving{existence part: the $c$-number lemma}.

\begin{remark}[on \cref{hyp:s03:T2}]\label{rem:s03:T2}
(i) This is a statement about lattice Feynman integrals: the difference of two mass counterterms at the same photon mass. The limit $\alpha_\infty(M)=\lim_k\alpha^\infty_k(M)$ and its description as a function of the couplings are discussed in Lemma~\ref{lem:s14:CW}.
(ii) \emph{Consumers.} Proposition~\ref{prop:s10:U} (the constants $\sigma_+=s_1$, $H$, $\kappa$ of the tube) and the coupling range, Lemma~\ref{lem:s14:CW}.
\end{remark}

\begin{hypothesis}[\Tthree: the threshold]\label{hyp:s03:T3}
$\sigma_1(M_1)>0$ for some $M_1\in(0,m)$.
\end{hypothesis}

\begin{remark}[on \Tthree]\label{rem:s03:T3}
(i) \Tthree\ is the only condition on the couplings anywhere in this paper, and it is one-sided: it fails exactly on the set $E^<_W$ of couplings $(e,\lambda,m,\lmu)$ at which $m$ minimises $M\mapsto\alpha_\infty(M)$ on $(0,m]$ (Lemma~\ref{lem:s14:CW}). At one loop it holds if and only if
\[
e^2<e_c^2(\lambda,m,\lmu)=\max\Bigl\{\frac{8\pi(m+\lmu)^2}{m+2\lmu}\Bigl(1+\frac{\lambda(N+2)}{2\pi m}\Bigr),\ 4\pi(m+\lmu)\Bigl(1+\frac{\lambda(N+2)}{\pi m}\Bigr)\Bigr\},
\]
$N$ the number of real components of the scalar field (equation~\eqref{eq:s14:ec2}); so at one loop a half-line of values of $e^2$ fails. Beyond one loop, \Tthree\ is \emph{proved} on explicit coupling regions (Proposition~\ref{prop:s14:T3}): at $\lmu=m$ on a region containing $0\le\lambda,e^2\le0.1m$, and at every $\lmu\in(0,m]$ on explicit boxes; the main theorem is stated on the region where it holds. The two-sided variant and the set $E_W\subseteq E^<_W$ are described in Lemma~\ref{lem:s14:CW}.
(ii) \emph{Consumers.} Proposition~\ref{prop:s10:U} and Lemma~\ref{lem:s14:CW}.
\end{remark}

\subsection{The rotation theorem}

For a family $f$ of real $2$-forms and $J\ge J_0(f)$ let $L^f_J$ be the set of locally uniform limits on $\C^N$ of subsequences of $(a^{(a)}_{k,J})_k$, $a^{(a)}_{k,J}(z)=\ava{e^{\sum_jz_jF(f_j)}}$ (the comparison family; by Lemma~\ref{lem:s05:C2} and Montel's theorem the set is non-empty).

\begin{theorem}[rotation invariance of the comparison family's limit states]\label{hyp:s03:RC}
For every finite family $f$ of real $2$-forms in $C^\infty_c(\R^3;\Lambda^2)$ and every choice $\zeta_J\in L^f_J$, $J\ge J_0(f)$, the limit $Z^{\mathrm{LIM}}_\infty(\sum_jz_jf_j,0)\defeq\lim_{J\to\infty}\zeta_J(z)$ exists locally uniformly on $\C^N$, is entire and does not depend on the choices. The functional $Z^{\mathrm{LIM}}_\infty(\cdot,0)$ is invariant under the Euclidean group $\Euc$: $Z^{\mathrm{LIM}}_\infty(\gamma\cdot f,0)=Z^{\mathrm{LIM}}_\infty(f,0)$ for every Euclidean motion $\gamma$ and every complex coefficient vector, $F$ transforming as a $2$-form.
\end{theorem}
Proved in the rotation part \moving{rotation part}: King's cubes are compared with periodic tori carrying a smooth mass wall, the walled tori on two lattices whose orientations differ by $\arctan(3/4)$ are compared, and King's state is identified as the lower envelope of the walled tori.

\begin{remark}[on \cref{hyp:s03:RC}]\label{rem:s03:RC}
(i) The theorem is posed for the comparison family and for the $J$-limit of its subsequential fixed-cube limits, not for the fixed-cube limits themselves; it does not give \cref{hyp:s03:EC}. Under \cref{hyp:s03:EC}, Corollary~\ref{cor:s04:G2} identifies $Z^{\mathrm{LIM}}_\infty(\cdot,0)$ with King's printed state $\ZKinf(\cdot,0)$, which is thereby $\Euc$-invariant (Theorem~\ref{thm:s06:O3}).
(ii) \emph{Consumers.} The $O(3)$ clause of the curvature sector (Theorem~\ref{thm:s06:O3}), the Euclidean clause of the measure's ergodic components (Theorem~\ref{thm:s09:S2K}(c)), and the $O(3)$ form of the LIM theorem (Theorem~\ref{thm:s12:LIM}).
\end{remark}

\begin{remark}[the Wick tube on King's convention]\label{rem:s03:TWK}
Theorem~\ref{thm:s11:W} is proved for the comparison convention $\NDa$. Its transfer to King's convention --- the $(k,J)$-uniform tube bound~\eqref{eq:s11:W} for $Z^{a,K}_{k,J}$ and the fixed-cube limits on the tube --- is Theorem~\ref{thm:s13:TWK}, with the kill step of Remark~\ref{rem:s04:F}, the cell statement of Lemma~\ref{lem:s04:E} and the direction check for the negative leg as its convention-sensitive steps; Form~(B) of the main theorem then holds on King's convention (Proposition~\ref{prop:s13:WK}).
\end{remark}

\subsection{Hypotheses of the optional extension}

The following statements are hypotheses of the optional extension of the Wick leg to arbitrary amplitudes (\cref{thm:s23:FormA}) only; no clause of \cref{thm:s13:main} uses them.

\begin{hypothesis}[\Wtwo: the Wick sector at arbitrary amplitudes]\label{hyp:s03:W2}
Let $a$ be an admissible scheme. (a) For every finite family of real $2$-forms and real functions, every $J\ge J_0(f,h)$ and \emph{every} real coefficient vector $(z,w)$, the limit $\lim_kZ^a_{k,J}(z,w)$ exists and lies in $(0,\infty)$. (b) For every $q\in C^\infty_c(\R^3)$ with $q\ge0$,
\[
\sup_{J\ge J_q}\ \lim_{k\to\infty}Z^{a}_{k,J}(0,q)<\infty ,
\]
where $J_q$ is the least $J$ with $\supp q$ at $\ell^\infty$-distance greater than $1$ from $\R^3\setminus\CJ$.
\end{hypothesis}

\begin{remark}[on \Wtwo]\label{rem:s03:W2}
(i) Clause (a) is \cref{hyp:s03:W1} with the threshold on the positive scalar part removed. On the cone of amplitudes $\sup q\le s_1<\sigma_1(M_1)$, clause (b) is proved in this paper in $L^1$ form (Corollary~\ref{cor:s11:W2prime}, on the comparison convention; on $\NDK$ in $\limsup$ form by Remark~\ref{rem:s04:F}). Above the threshold both clauses are open; the exact obstruction is the free-energy bound $\Bminus{u}$ at amplitude $u$ above the threshold (Hypothesis~\ref{hyp:s23:Bminus}): the $\eps$-uniform stability of King's torus model with its bare coefficient shifted by $-2u$.
(ii) \emph{Consumers.} None in the headline. The headline Wick result is Form~(B) on the tube $\Dtube$, where only the proved cone instance enters, through Theorem~\ref{thm:s11:W}; \Wtwo\ is a hypothesis of Form~(A) (Theorem~\ref{thm:s23:FormA}) and of \RW.
\end{remark}

\begin{hypothesis}[\RW: rotation invariance of the Wick sector at arbitrary amplitudes]\label{hyp:s03:RW}
For an admissible scheme $a$ satisfying \Wtwo, the full object $Z^a_\infty=\lim_J\lim_kZ^a_{k,J}$ of the comparison family (whose existence and entireness follow from \Wtwo\ as in Lemma~\ref{lem:s08:W1}) is invariant under every Euclidean motion for all complex coefficients, $F$ transforming as a $2$-form and $W^a$ as a scalar.
\end{hypothesis}

\begin{remark}[on \RW]\label{rem:s03:RW}
\RW\ enters only in the Euclidean clause of Form~(A) (Theorem~\ref{thm:s23:FormA}). The headline Wick result, Form~(B) on the tube, has no Euclidean clause and does not assume it. The rotation part proves the curvature sector only.
\end{remark}

\begin{hypothesis}[\FAK: Form~(A) on King's convention]\label{hyp:s03:FAK}
On King's convention $\NDK$, under \Wtwo\ read for $\NDK$: the limit functional $Z^{a,K}_\infty$ is entire in all coefficients; the lattice symmetries $\Oh$ about every point hold for the joint field at all amplitudes; reflection positivity holds at arbitrary complex coefficients; the measure exists off the tube; and the Euclidean clause of the Wick sector holds for the King-family object.
\end{hypothesis}

\begin{remark}[on \FAK]\label{rem:s03:FAK}
\FAK\ is the King-convention counterpart of Form~(A) as the comparison-convention results give it. It is consumed only in the optional extension (Theorem~\ref{thm:s23:FormA}) and is listed in the open problems; the headline makes no use of it.
\end{remark}

\subsection{Hypotheses stated in later sections}

The statements on continuum subsequences (Theorem~\ref{thm:s12:LIM}) are proved under three hypotheses stated there, \Bcell\ (Hypothesis~\ref{hyp:s12:Bcell}; the cell bound, an instance of \cref{hyp:s03:W1}), \SCH\ (Hypothesis~\ref{hyp:s12:SCH}; $a=z$ or $\sup_k\abs{z_k-a_k}<\infty$) and \UNIQ\ (Hypothesis~\ref{hyp:s12:UNIQ}; any two continuum subsequences give the same LIM state). They are hypotheses of those statements only: all three are proved for the objects of the main theorem (Proposition~\ref{prop:s12:ID}), and none is a hypothesis of Theorem~\ref{thm:s13:main}. The optional extension (Theorem~\ref{thm:s23:FormA}) uses $\Bminus{u}$ (Hypothesis~\ref{hyp:s23:Bminus}), the obstruction behind \Wtwo\ above the threshold. The limit of vanishing photon mass is not treated in this paper.

\subsection{Results of the earlier paper used as inputs}

\begin{remark}[results of the torus paper used as proved]\label{rem:s03:companion-inputs}
The following results of the authors' earlier paper on the torus limits \cite{SymPhase} are used as proved and cited by their labels there; none is re-proved here, and each is used only under the hypotheses it carries there, which are matched to King's state at the places where they are applied (Theorem~\ref{thm:s09:S2K}, Proposition~\ref{prop:s14:T3}): its abstract chain from a limit functional with the properties of Theorem~\ref{thm:s05:C} and the reflection positivity of Theorem~\ref{thm:s08:OS2C} to the measure, its ergodic decomposition and weak clustering --- thm:D:decomp, lem:D:realize, lem:D:hermite, prop:discharge, thm:D:mean, thm:D:diag, prop:D:recon, prop:D:fixed(iii), cor:D:schwinger(iv), add:Ohdecomp, add:OS4, p:cor:rotcor, prop:C-tight, prop:C-RP(a), lem:C-FP and the lemmas of its Appendix~G1 (the two-loop majorants behind Proposition~\ref{prop:s14:T3}); Theorem~main(ii); the lemmas on the flat state and the current representation, lem:flat, lem:C-flatlimit, lem:C-current, lem:C-Fsector.
Of \cite{BFS1}, Theorems~6.1 and~6.2(1),(2) are cited as printed; items (3),(4) of its Theorem~6.2, a sketch in print, are a lemma of this paper (Lemma~\ref{lem:s02:BFS}). The chessboard estimate is Theorem~4.1/4.2 of \cite{FILS78}, applied after its hypotheses are checked (Proposition~\ref{prop:s10:RP}). Nothing of \cite{K1,K2} is used as a result: King's existence theorem is proved in \cref{sec:k01} and \moving{existence section with scalar sources}, and his reflection-positivity argument for the modified box and his passage to subsequences are replaced by proofs in this paper.
\end{remark}
